% SUPERSEDED 2026-10-28 ~19:15 UTC: moved into results.tex (sec:res-sense, with Figure fig:tilt-sense); edit there.
% DRAFT (2026-10-28 evening, not \input anywhere): "why the cation sets the tilt sense", to go after the Table 8
% paragraph of sec:res-tilt once the oriented-cation DFT (tilt_sense/orient) and the literature check
% (tilt_sense/literature_tilt_sense.md) are in. Sources: tilt_sense/tilt_scan (build_tilt_scan.py, tilt_scan_relax.py,
% tilt_scan_analyze.py, tilt_scan_orient.py), tilt_sense/tilt_sense_compare.md (48-atom DFT).
\subsection{\protect\changed{Why the cation sets the tilt sense}}
\label{sec:res-sense}
{\changedcolor
The opposite stacking of FA and MA compounds is well documented for the lead iodides and their
mixtures.\cite{Fabini2016Angew, Hainer2025NatCommun, Dubajic2025NatNano} The inorganic framework hardly prefers
either: in DFT the in-phase and antiphase tilts of \CsSnI{} at 8 and \ang{12} differ by at most
\SI{0.04}{\milli\hartree} per formula unit with either functional, as found for CsPbI$_3$ and
\CsSnI{},\cite{Klarbring2019PRB} against 0.4 to \SI{2.8}{\milli\hartree} for the hybrids. The sense is chosen by the
A-site cation. Proposed reasons are a general preference of halides for in-phase tilts,\cite{Young2016JPCL} hydrogen
bonds of MA that drive the antiphase mode,\cite{Lee2016ChemMater} an FA arrangement that templates the in-phase
tilt,\cite{Senno2021PCCP} and the near-degeneracy of the two tilt modes.\cite{Hainer2025NatCommun} We show that the
shape of the cage accounts for both cations.

The static energy of a single cell does not show which sense a cation prefers, because it depends on how the cations
are arranged. In the four-formula-unit cells used for DFT, the antiphase--in-phase difference changes by 2 to
\SI{3}{\milli\hartree} per formula unit from one random arrangement to the next, more than its mean. We therefore
computed it with each potential for 20 to 160 random arrangements in cells of 4 to 128 formula units, tilted in-phase
or antiphase by 4, 8 or \ang{12} about $c$ at a fixed cubic metric. Each arrangement is placed in both cells and the
cations are relaxed with Sn and I fixed. The scatter falls as the inverse square root of the cell size, and the
largest cells give a clear answer (Figure~\ref{fig:tilt-sense}). In the 128-formula-unit cells at \ang{8},
\FASnI{} is lower in-phase by \num{0.88 \pm 0.08} (PBEsol) and \SI{1.19 \pm 0.12}{\milli\hartree} (\rSCAN{}) per
formula unit for all but one arrangement, and \MASnI{} is lower antiphase by \num{0.62 \pm 0.05} and
\SI{1.23 \pm 0.08}{\milli\hartree} for every arrangement. Both preferences peak near \ang{8} and fall to a few tenths of a
millihartree at \ang{12}, where \FASnI{} turns slightly antiphase. The sense is therefore chosen where the tilt forms,
near the transition; at large static tilts FAPbI$_3$ is antiphase as well.\cite{Dutta2025JACS}

The cause is the shape of the cage. A tilt about $c$ turns each face of the A-site cage into a rectangle, long along
one in-plane axis $L$. In-phase tilting gives the upper and lower faces the same long axis, a box of uniform
cross-section; antiphase tilting crosses them, so the long axis turns by \ang{90} from the lower to the upper face.
Fitting the energy difference of the small cells to the cation orientations gives the contribution of a cation
pointing along each axis. Both cations penalise in-phase tilting when they lie along the short axis of the box and
favour it when they lie along $L$. They differ along $c$: an FA cation standing along $c$ favours in-phase by 6 to
\SI{8}{\milli\hartree}, an MA cation favours antiphase by 3 to \SI{4}{\milli\hartree}. Planar FA carries an NH$_2$
group at each end; lying along $L$, or standing along $c$ with its plane through $L$, it meets two equivalent faces of
the box, while in the crossed cage one end is always squeezed. Averaged over orientations, FA favours in-phase and MA
antiphase. Relaxation turns the FA cations towards $L$ in the in-phase cell ($\langle (\hat u\cdot L)^2\rangle$ rises
from 0.33 to 0.6--0.7 for the N--N axis $\hat u$).

The deposited structures and the MD agree. In both structures of \FASnI{} (Table~\ref{tab:expt-tilt}) the FA
nitrogens lie along the long axis of their cage, with the N--N axes of neighbouring cages perpendicular in the plane and
repeated along $c$;\cite{Kahmann2020ACSEL} the refinement of tetragonal \MASnI{} places the C--N axis along
$c$,\cite{stoumpos2013semiconducting} the orientation that favours antiphase stacking most, although with large
displacement parameters. In the in-phase phase of the MD the FA cations keep their N--N axis in the plane (mean
$|\cos|$ with the tilt axis 0.16 to 0.29, against 0.5 for random orientations), while in the antiphase phase the MA
cations turn almost freely (0.46 to 0.48), so for MA the preference is an average over orientations.
% TODO(oriented DFT): one sentence with the DFT energies of the oriented cells (FA Lflat / Lpflat / cL / cLp; MA L / Lp / c).
% literature: done (JCTC.bib keys Klarbring2019PRB, Young2016JPCL, Lee2016ChemMater, Senno2021PCCP, Hainer2025NatCommun,
% Dubajic2025NatNano, Dutta2025JACS, Fabini2016Angew; tilt_sense/literature_tilt_sense.md).
% MA/FA orientation: done (cage_sense.py on x5_2304_<MA|FA>_R2_cool_r05; CIF check of COD 4517852/4517853/4335636 with
% tilt_scan_orient.long_axes: (N-N.L)^2 = 1.00 in both FASnI3 structures; MASnI3 I4cm C-N along c, Uiso 0.17/0.27).
}
